\textbf{Keywords:} WSL, Git, AICEX, ngspice, xschem, Magic, cicconf,
cicsim, cicpy

\section{Tools}\label{tools}\index{tools@Tools}

I would strongly recommend that you install all tools locally on your
system. There is a video that describes the install procedure. It's a
few years old, but should still be able to guide you
\url{https://youtu.be/DRppsdjo2Rc?si=x8cJsa1lpncvSFmu}.

For the analog toolchain we need some tools, and a process design kit
(PDK).

\begin{itemize}
\tightlist
\item
  \href{https://github.com/google/skywater-pdk}{Skywater 130nm PDK}. I
  use \href{https://github.com/RTimothyEdwards/open_pdks}{open\_pdks} to
  install the PDK
\item
  \href{https://github.com/RTimothyEdwards/magic}{Magic VLSI} for layout
  (Version 8.3 revision 541)
\item
  \href{https://git.code.sf.net/p/ngspice/ngspice}{ngspice} for
  simulation (version 45.2)
\item
  \href{https://github.com/RTimothyEdwards/netgen.git}{netgen} for LVS
  (1.5.295)
\item
  \href{https://github.com/StefanSchippers/xschem}{xschem} (3.4.8RC)
\item
  \href{https://www.veripool.org/verilator/}{verilator} (5.034)
\item
  python \textgreater{} 3.10
\end{itemize}

The tools are not that big, but the PDK is huge, so you need to have
about 50 GB disk space available.

\subsection{Setup WSL (Applicable for Windows
users)}\label{setup-wsl-applicable-for-windows-users}

Install a Linux distribution such as Ubuntu 24.04 LTS by running the
following command in PowerShell on Windows and follow the instructions.

\begin{Shaded}
\begin{Highlighting}[]
\ExtensionTok{wsl} \AttributeTok{{-}{-}install} \AttributeTok{{-}d}\NormalTok{ Ubuntu{-}24.04}
\end{Highlighting}
\end{Shaded}

When you have installed the Linux distribution and signed into it,
install make

\begin{Shaded}
\begin{Highlighting}[]
\FunctionTok{sudo}\NormalTok{ apt install make}
\end{Highlighting}
\end{Shaded}

\subsection{Setup public key towards
github}\label{setup-public-key-towards-github}

Do

\begin{Shaded}
\begin{Highlighting}[]
\FunctionTok{ssh{-}keygen} \AttributeTok{{-}t}\NormalTok{ rsa}
\end{Highlighting}
\end{Shaded}

And press ``enter'' on most things, or if you're paranoid, add a
passphrase

Then

\begin{Shaded}
\begin{Highlighting}[]
\FunctionTok{cat}\NormalTok{ \textasciitilde{}/.ssh/id\_rsa.pub }
\end{Highlighting}
\end{Shaded}

And add the public key to your github account. Settings - SSH and GPG
keys

\subsection{Provide git with author
identity}\label{provide-git-with-author-identity}

There are interactions with git that require an author identity. You are
supposed to use one of these interactions a lot during the project,
namely, \texttt{git\ commit}. What you need to provide is an email
address and a name. If you would like to keep your real email address
private/secret, read what it says on GitHub at your user settings page
under \href{https://github.com/settings/emails}{emails}. Use the below
commands to provide the author identity information to git.

\begin{Shaded}
\begin{Highlighting}[]
\FunctionTok{git}\NormalTok{ config }\AttributeTok{{-}{-}global}\NormalTok{ user.email }\StringTok{"you@example.com"}
\FunctionTok{git}\NormalTok{ config }\AttributeTok{{-}{-}global}\NormalTok{ user.name }\StringTok{"Your Name"}
\end{Highlighting}
\end{Shaded}

\subsection{Get AICEX and setup your
shell}\label{get-aicex-and-setup-your-shell}

You don't have to put aicex in \texttt{\$HOME/pro}, but if you don't
know where to put it, choose that directory.

\begin{Shaded}
\begin{Highlighting}[]
\BuiltInTok{cd} 
\FunctionTok{mkdir}\NormalTok{ pro }
\BuiltInTok{cd}\NormalTok{ pro}
\FunctionTok{git}\NormalTok{ clone }\AttributeTok{{-}{-}recursive}\NormalTok{ https://github.com/wulffern/aicex.git}
\end{Highlighting}
\end{Shaded}

You need to add the following to your \texttt{\textasciitilde{}/.bashrc}
(note that \texttt{\textasciitilde{}} refers to your home directory
\texttt{\$HOME/.bashrc} also works, or \texttt{\$HOME/.bash\_profile} on
some newer macs)

\begin{Shaded}
\begin{Highlighting}[]
\BuiltInTok{export} \VariableTok{PDK\_ROOT}\OperatorTok{=}\NormalTok{/opt/pdk/share/pdk}
\BuiltInTok{export} \VariableTok{LD\_LIBRARY\_PATH}\OperatorTok{=}\NormalTok{/opt/eda/lib}
\BuiltInTok{export} \VariableTok{PATH}\OperatorTok{=}\NormalTok{/opt/eda/bin:}\VariableTok{$HOME}\NormalTok{/.local/bin:}\VariableTok{$PATH}
\end{Highlighting}
\end{Shaded}

\subsection{On systems with python3 \textgreater{}
3.12}\label{on-systems-with-python3-3.12}

On newer systems it's not trivial to install python packages because
python is externally managed. As such, we need to install a python
environment.

\begin{Shaded}
\begin{Highlighting}[]
\CommentTok{\#{-} Find a package similar to name below}
\FunctionTok{sudo}\NormalTok{ apt{-}get update}
\FunctionTok{sudo}\NormalTok{ apt install python3.12{-}venv}
\FunctionTok{sudo}\NormalTok{ mkdir /opt}
\FunctionTok{sudo}\NormalTok{ mkdir /opt/eda}
\FunctionTok{sudo}\NormalTok{ mkdir /opt/eda/python3}
\FunctionTok{sudo}\NormalTok{ chown }\AttributeTok{{-}R} \VariableTok{$USER}\NormalTok{:}\VariableTok{$USER}\NormalTok{ /opt/eda/python3/}
\ExtensionTok{python3} \AttributeTok{{-}m}\NormalTok{ venv /opt/eda/python3}
\end{Highlighting}
\end{Shaded}

Modify the \texttt{\textasciitilde{}/.bashrc} to include the python
environment

\begin{Shaded}
\begin{Highlighting}[]
\BuiltInTok{export} \VariableTok{PATH}\OperatorTok{=}\NormalTok{/opt/eda/bin:/opt/eda/python3/bin:}\VariableTok{$HOME}\NormalTok{/.local/bin:}\VariableTok{$PATH}
\end{Highlighting}
\end{Shaded}

\subsection{Install Tools}\label{install-tools}\index{install tools@Install tools}

Make sure you load the settings before you proceed

\begin{Shaded}
\begin{Highlighting}[]
\BuiltInTok{source}\NormalTok{ \textasciitilde{}/.bashrc}
\end{Highlighting}
\end{Shaded}

Hopefully the commands below work, if not, then try again, or try to
understand what fails. There is no point in continuing if one command
fails.

\begin{Shaded}
\begin{Highlighting}[]
\BuiltInTok{cd}\NormalTok{ aicex/tests/}
\FunctionTok{make}\NormalTok{ requirements}
\FunctionTok{make}\NormalTok{ tt}
\end{Highlighting}
\end{Shaded}

On a mac, you probably need to add bison to the path

\begin{Shaded}
\begin{Highlighting}[]
\BuiltInTok{export} \VariableTok{PATH}\OperatorTok{=}\StringTok{"/opt/homebrew/opt/bison/bin:}\VariableTok{$PATH}\StringTok{"}
\end{Highlighting}
\end{Shaded}

I've split the install of each of the tools. It's possible to run the
commented out lines instead, but they often fail

\begin{Shaded}
\begin{Highlighting}[]
\CommentTok{\#make eda\_compile}
\CommentTok{\#sudo make eda\_install}
\FunctionTok{make}\NormalTok{ magic\_compile  magic\_install}
\FunctionTok{make}\NormalTok{ netgen\_compile netgen\_install}
\FunctionTok{make}\NormalTok{ xschem\_compile xschem\_install}
\FunctionTok{make}\NormalTok{ iverilog\_compile iverilog\_install}
\FunctionTok{make}\NormalTok{ ngspice\_compile }\CommentTok{\# Sometimes fails}
\FunctionTok{make}\NormalTok{ ngspice\_compile ngspice\_install}
\end{Highlighting}
\end{Shaded}

For later versions of ngspice, they have modified the build system, so
you may have to run

\begin{verbatim}
make ngspice2_compile
make ngspice2_install
\end{verbatim}

On Mac, do

\begin{Shaded}
\begin{Highlighting}[]
\ExtensionTok{brew}\NormalTok{ install yosys verilator}
\end{Highlighting}
\end{Shaded}

On Linux, do

\begin{Shaded}
\begin{Highlighting}[]
\FunctionTok{make}\NormalTok{ yosys\_compile yosys\_install}
\end{Highlighting}
\end{Shaded}

On all, do

\begin{Shaded}
\begin{Highlighting}[]
\ExtensionTok{python3} \AttributeTok{{-}m}\NormalTok{ ensurepip }\AttributeTok{{-}{-}default{-}pip}

\ExtensionTok{python3} \AttributeTok{{-}m}\NormalTok{ pip install matplotlib numpy click svgwrite }\DataTypeTok{\textbackslash{}}
\NormalTok{    pyyaml pandas tabulate wheel setuptools tikzplotlib}
\BuiltInTok{source}\NormalTok{ install\_open\_pdk.sh}
\end{Highlighting}
\end{Shaded}

\subsection{Install cicconf}\label{install-cicconf}\index{install cicconf@Install cicconf}

cIcConf is used for configuration. How the IPs are connected, and what
version of IPs to get.

\begin{Shaded}
\begin{Highlighting}[]
\BuiltInTok{cd}
\BuiltInTok{cd}\NormalTok{ pro/aicex/ip/cicconf}
\FunctionTok{git}\NormalTok{ checkout main }
\FunctionTok{git}\NormalTok{ pull}
\ExtensionTok{python3} \AttributeTok{{-}m}\NormalTok{ pip install }\AttributeTok{{-}e}\NormalTok{ .}
\BuiltInTok{cd}\NormalTok{ ../}
\end{Highlighting}
\end{Shaded}

Update IPs

\begin{Shaded}
\begin{Highlighting}[]
\ExtensionTok{cicconf}\NormalTok{ clone }\AttributeTok{{-}{-}https} 
\BuiltInTok{cd}\NormalTok{ ../..}
\end{Highlighting}
\end{Shaded}

\subsection{Install cicsim}\label{install-cicsim}\index{install cicsim@Install cicsim}

cIcSim is used for simulation orchestration.

\begin{Shaded}
\begin{Highlighting}[]
\BuiltInTok{cd}\NormalTok{ aicex/ip/cicsim}
\FunctionTok{git}\NormalTok{ checkout main}
\FunctionTok{git}\NormalTok{ pull}
\ExtensionTok{python3} \AttributeTok{{-}m}\NormalTok{ pip install }\AttributeTok{{-}e}\NormalTok{ .}
\BuiltInTok{cd}\NormalTok{ ../..}
\end{Highlighting}
\end{Shaded}

\subsection{Install cicpy}\label{install-cicpy}\index{install cicpy@Install cicpy}

CicPy is used to generate layout

\begin{Shaded}
\begin{Highlighting}[]
\BuiltInTok{cd}\NormalTok{ aicex/ip/cicpy }
\FunctionTok{git}\NormalTok{ checkout master}
\FunctionTok{git}\NormalTok{ pull}
\ExtensionTok{python3} \AttributeTok{{-}m}\NormalTok{ pip install }\AttributeTok{{-}e}\NormalTok{ .}
\BuiltInTok{cd}\NormalTok{ ../..}
\BuiltInTok{cd}\NormalTok{ aicex/ip/cicspi }
\FunctionTok{git}\NormalTok{ checkout main }
\FunctionTok{git}\NormalTok{ pull}
\ExtensionTok{python3} \AttributeTok{{-}m}\NormalTok{ pip install }\AttributeTok{{-}e}\NormalTok{ .}
\BuiltInTok{cd}\NormalTok{ ../..}
\end{Highlighting}
\end{Shaded}

\subsection{Setup your ngspice
settings}\label{setup-your-ngspice-settings}

Edit \texttt{\textasciitilde{}/.spiceinit} and add

\begin{Shaded}
\begin{Highlighting}[]
\BuiltInTok{set}\NormalTok{ ngbehavior=hsa     }\KeywordTok{;} \BuiltInTok{set}\NormalTok{ compatibility for PDK libs}
\BuiltInTok{set}\NormalTok{ ng\_nomodcheck      }\KeywordTok{;} \ExtensionTok{don}\StringTok{\textquotesingle{}t check the model parameters}
\StringTok{set num\_threads=8      ; CPU hardware threads available}
\StringTok{set skywaterpdk}
\StringTok{option noinit          ; don\textquotesingle{}}\ExtensionTok{t}\NormalTok{ print operating point data}
\ExtensionTok{option}\NormalTok{ klu}
\ExtensionTok{optran}\NormalTok{ 0 0 0 100p 2n 0 }\KeywordTok{;} \ExtensionTok{don}\StringTok{\textquotesingle{}t use dc operating point,}
\StringTok{option opts}
\end{Highlighting}
\end{Shaded}

\section{Check that magic and xschem
work}\label{check-that-magic-and-xschem-work}

To check that magic and xschem work

\begin{Shaded}
\begin{Highlighting}[]
\BuiltInTok{cd}\NormalTok{ \textasciitilde{}/pro/aicex/ip/sun\_sar9b\_sky130nm/work }
\ExtensionTok{magic}\NormalTok{ ../design/SUN\_SAR9B\_SKY130NM/SUNSAR\_SAR9B\_CV.mag }\KeywordTok{\&}
\ExtensionTok{xschem} \AttributeTok{{-}b}\NormalTok{ ../design/SUN\_SAR9B\_SKY130NM/SUNSAR\_SAR9B\_CV.sch }\KeywordTok{\&}
\end{Highlighting}
\end{Shaded}

\section{Summary}\label{summary}

The one-page version of this chapter:

\begin{itemize}
\tightlist
\item
  The whole flow is open source: xschem, ngspice, magic, netgen, and the
  cic tools on top
\item
  Everything installs from scripts, and the docs and video walk the
  setup end to end
\item
  If the tools do not run, nothing else in this course happens - do this
  first, and ask early when stuck
\end{itemize}

\section{Would you like to know
more?}\label{would-you-like-to-know-more}

Each tool documents itself well: xschem, ngspice, magic, netgen and the
SKY130 PDK

The tools chapter later in this book covers the cic layer built on top
of them
